\chapter{پیچیدگی زمانی و کلاس P}

حتی وقتی مسئله‌ای تصمیم‌پذیر است، ممکن است در عمل حل‌شدنی نباشد، چون الگوریتم آن زمان یا حافظه بیش از حدی لازم دارد. \textbf{نظریه پیچیدگی محاسبه} منابع لازم برای حل مسئله‌ها (زمان، حافظه، دیسک، تعداد پیام و \dots) را بررسی می‌کند. در این بخش روی \textbf{زمان} تمرکز می‌کنیم: یاد می‌گیریم زمان را چطور اندازه بگیریم و مسئله‌ها را بر اساس زمان لازم دسته‌بندی کنیم.

\section{اندازه‌گیری زمان}

\begin{example}
زبان $A=\{0^k1^k\st k\ge0\}$ تصمیم‌پذیر است. یک ماشین تورینگ تک‌نواره برای تصمیم‌گیری آن چقدر زمان لازم دارد؟
\begin{algo}{$M_1$}{روی ورودی $w$}
\begin{enumerate}
\item نوار را طی کن و اگر یک $0$ در سمت راست یک $1$ پیدا شد، \textbf{رد کن}.
\item تا وقتی هم $0$ و هم $1$ روی نوار مانده، تکرار کن:
\item \quad نوار را طی کن و یک $0$ و یک $1$ را خط بزن.
\item اگر بعد از خط خوردن همه $1$ها هنوز $0$ مانده بود یا برعکس، \textbf{رد کن}. وگرنه، اگر نه $0$ و نه $1$ روی نوار نمانده، \textbf{بپذیر}.
\end{enumerate}
\end{algo}
\end{example}

زمان اجرا به عوامل زیادی بستگی دارد. برای ساده شدن، فقط آن را تابعی از \textbf{طول نمایش ورودی} در نظر می‌گیریم. همچنین معمولاً \textbf{تحلیل بدترین حالت} انجام می‌دهیم: بیشترین زمان اجرا روی همه ورودی‌های به طول $n$. (در \textbf{تحلیل حالت متوسط}، میانگین زمان روی همه ورودی‌های به طول $n$ حساب می‌شود.)

\begin{definition}[زمان اجرا]
فرض کنید $M$ یک ماشین تورینگ قطعی باشد که روی همه ورودی‌ها متوقف می‌شود. \textbf{زمان اجرا} یا \textbf{پیچیدگی زمانی} $M$ تابع $f:\mathbb{N}\to\mathbb{N}$ است که $f(n)$ بیشترین تعداد گامی است که $M$ روی ورودی‌های به طول $n$ برمی‌دارد. می‌گوییم $M$ در زمان $f(n)$ اجرا می‌شود یا $M$ یک ماشین با زمان $f(n)$ است. معمولاً $n$ طول ورودی است.
\end{definition}

\subsection{نمادهای $O$ بزرگ و $o$ کوچک}

زمان دقیق اجرا معمولاً عبارت پیچیده‌ای است. پس آن را \textbf{تخمین} می‌زنیم: با \textbf{تحلیل مجانبی} رفتار روی ورودی‌های بزرگ را بررسی می‌کنیم. فقط جمله با بالاترین مرتبه را نگه می‌داریم و ضریب آن و جمله‌های مرتبه پایین‌تر را کنار می‌گذاریم. مثلاً
\[f(n)=6n^3+2n^2+20n+45\quad\leadsto\quad n^3,\]
یعنی $f$ به‌طور مجانبی حداکثر $n^3$ است و می‌نویسیم $f(n)=O(n^3)$.

\begin{definition}[$O$ بزرگ]
فرض کنید $f,g:\mathbb{N}\to\mathbb{R}^+$. می‌گوییم $f(n)=O(g(n))$ اگر اعداد صحیح مثبت $c$ و $n_0$ وجود داشته باشند که برای هر $n\ge n_0$:
\[f(n)\le c\,g(n).\]
در این حالت $g(n)$ یک \textbf{کران بالا} (یا دقیق‌تر، کران بالای مجانبی) برای $f(n)$ است.
\end{definition}

یعنی اگر تفاوت در حد ضریب ثابت را نادیده بگیریم، $f$ کوچک‌تر یا مساوی $g$ است؛ $O$ نماینده یک ثابت پنهان است.

\begin{example}
\begin{itemize}
\item $f_1(n)=5n^3+2n^2+22n+6=O(n^3)$، چون برای $n\ge10$ داریم $f_1(n)\le6n^3$. همچنین $f_1(n)=O(n^4)$، اما $f_1(n)$ از مرتبه $O(n^2)$ \textbf{نیست}.
\item تغییر پایه لگاریتم مقدار را فقط در یک ضریب ثابت تغییر می‌دهد: $\log_bn=\log_2n/\log_2b$. پس در نماد $O$ پایه را نمی‌نویسیم: $f(n)=O(\log n)$.
\item نماد $O$ را در عبارت‌های حسابی هم به کار می‌بریم: $f(n)=O(n^2)+O(n)=O(n^2)$.
\item $O$ در توان: $f(n)=2^{O(n)}$ یعنی $f(n)\le2^{cn}$ برای یک $c$. همچنین $f(n)=2^{O(\log n)}$ یعنی $f(n)\le n^c$ برای یک $c$ (چون $n=2^{\log_2n}$)، و $f(n)=n^{O(1)}$ هم یعنی $f(n)\le n^c$.
\item کران‌های \textbf{چندجمله‌ای} شکل $n^c$ دارند ($c>0$) و کران‌های \textbf{نمایی} شکل $2^{n^\delta}$ ($\delta>0$).
\end{itemize}
\end{example}

برای اینکه بگوییم یک تابع به‌طور مجانبی \emph{اکیداً} کوچک‌تر از دیگری است، از $o$ کوچک استفاده می‌کنیم. رابطه $O$ و $o$ مثل رابطه $\le$ و $<$ است.

\begin{definition}[$o$ کوچک]
فرض کنید $f,g:\mathbb{N}\to\mathbb{R}^+$. می‌گوییم $f(n)=o(g(n))$ اگر
\[\lim_{n\to\infty}\frac{f(n)}{g(n)}=0.\]
به بیان دیگر برای هر عدد حقیقی $c>0$ عددی مثل $n_0$ وجود دارد که برای هر $n\ge n_0$: $f(n)<c\,g(n)$.
\end{definition}

مثلاً $n^2=o(n^3)$، اما $f(n)$ هرگز $o(f(n))$ نیست.

\subsection{تحلیل الگوریتم}

\begin{example}[تحلیل $M_1$]
\begin{itemize}
\item \textbf{مرحله ۱:} یک بار طی کردن نوار و برگشتن: $2n=O(n)$ گام.
\item \textbf{مراحل ۲ و ۳:} هر بار دو نماد خط می‌خورد، پس حداکثر $n/2$ دور داریم و هر دور $O(n)$ گام دارد: $(n/2)\,O(n)=O(n^2)$.
\item \textbf{مرحله ۴:} یک بار طی کردن نوار: $O(n)$.
\end{itemize}
پس زمان اجرای $M_1$ برابر $O(n)+O(n^2)+O(n)=O(n^2)$ است.
\end{example}

\begin{definition}[کلاس پیچیدگی زمانی]
فرض کنید $t:\mathbb{N}\to\mathbb{R}^+$ یک تابع باشد. \textbf{کلاس پیچیدگی زمانی} $\mathrm{TIME}(t(n))$ مجموعه همه زبان‌هایی است که یک ماشین تورینگ با زمان $O(t(n))$ آن‌ها را تصمیم می‌گیرد.
\end{definition}

پس $A=\{0^k1^k\st k\ge0\}\in\mathrm{TIME}(n^2)$. آیا می‌توان $A$ را سریع‌تر تصمیم گرفت؟

\begin{example}[الگوریتم $O(n\log n)$]
\begin{algo}{$M_2$}{روی ورودی $w$}
\begin{enumerate}
\item نوار را طی کن و اگر یک $0$ در سمت راست یک $1$ پیدا شد، \textbf{رد کن}.
\item تا وقتی هم $0$ و هم $1$ روی نوار مانده، تکرار کن:
\item \quad نوار را طی کن و بررسی کن که تعداد کل $0$ها و $1$های باقی‌مانده زوج است یا فرد. اگر فرد بود، \textbf{رد کن}.
\item \quad دوباره نوار را طی کن و یکی در میان $0$ها (با شروع از اولین $0$) و یکی در میان $1$ها (با شروع از اولین $1$) را خط بزن.
\item اگر هیچ $0$ و هیچ $1$ی روی نوار نمانده، \textbf{بپذیر}؛ وگرنه \textbf{رد کن}.
\end{enumerate}
\end{algo}
هر دور، تعداد $0$ها و تعداد $1$ها را نصف می‌کند، پس حداکثر $1+\log_2n$ دور داریم و هر دور $O(n)$ گام دارد. کل زمان $O(n\log n)$ است. (درستی: مرحله ۳ در هر دور، زوج یا فرد بودن تعداد $0$ها و $1$ها را مقایسه می‌کند. پشت سر هم قرار گرفتن این نتیجه‌ها در دورهای متوالی، دقیقاً نمایش دودویی تعداد $0$ها و تعداد $1$ها را از کم‌ارزش‌ترین بیت مقایسه می‌کند.)
\end{example}

\begin{example}[با ماشین دونواره: $O(n)$]
\begin{algo}{$M_3$}{روی ورودی $w$ (ماشین دونواره)}
\begin{enumerate}
\item نوار ۱ را طی کن و اگر یک $0$ در سمت راست یک $1$ پیدا شد، \textbf{رد کن}.
\item $0$های نوار ۱ را تا اولین $1$ طی کن و هم‌زمان آن‌ها را روی نوار ۲ کپی کن.
\item $1$های نوار ۱ را تا انتهای ورودی طی کن؛ به ازای هر $1$ یک $0$ را روی نوار ۲ خط بزن. اگر همه $0$ها قبل از تمام شدن $1$ها خط خوردند، \textbf{رد کن}.
\item اگر همه $0$ها خط خورده‌اند \textbf{بپذیر}؛ اگر $0$ی مانده \textbf{رد کن}.
\end{enumerate}
\end{algo}
هر مرحله $O(n)$ گام دارد، پس زمان کل $O(n)$ است.
\end{example}

\section{رابطه پیچیدگی مدل‌های محاسباتی}

طبق تز چرچ–تورینگ همه مدل‌های معقول محاسبه از نظر \emph{تصمیم‌پذیری} هم‌ارزند. اما در نظریه پیچیدگی، انتخاب مدل روی پیچیدگی زمانی زبان اثر دارد؛ مثلاً $A$ را با ماشین تک‌نواره در $O(n\log n)$ و با ماشین دونواره در $O(n)$ تصمیم گرفتیم. سه مدل را مقایسه می‌کنیم: ماشین تک‌نواره، چندنواره و غیرقطعی.

\begin{definition}[زمان اجرای ماشین غیرقطعی]
فرض کنید $N$ یک ماشین تورینگ غیرقطعی و تصمیم‌گیرنده باشد. \textbf{زمان اجرای} $N$ تابع $f:\mathbb{N}\to\mathbb{N}$ است که $f(n)$ بیشترین تعداد گامی است که $N$ روی \textbf{هر شاخه} از محاسبه‌اش روی ورودی‌های به طول $n$ برمی‌دارد.
\end{definition}
\begin{ltrfig}
\begin{tikzpicture}[font=\small,every node/.style={circle,fill=jzMain,inner sep=1.2pt}]
  \node (a0) at (0,0) {}; \foreach \i in {1,...,4} {\node (a\i) at (0,-\i*0.6) {};} \draw[jzMain] (a0)--(a1)--(a2)--(a3)--(a4);
  \node[fill=none,right] at (a4) {accept/reject}; \node[fill=none] at (0,0.5) {Deterministic};
  \draw[<->] (-0.8,0) -- node[fill=none,left] {$f(n)$} (-0.8,-2.4);
  \begin{scope}[xshift=5.5cm,edge from parent/.style={draw,jzMain},level distance=6mm,sibling distance=7mm]
    \node {} child {node {} child {node {} child {node {}}} child {node {}}} child {node {} child {node {} child {node {} child {node {}}}}} child {node {} child {node {}}};
    \node[fill=none] at (0,0.5) {Nondeterministic};
    \draw[<->] (1.8,0) -- node[fill=none,right] {$f(n)$} (1.8,-2.4);
  \end{scope}
\end{tikzpicture}
\end{ltrfig}

\begin{theorem}
فرض کنید $t(n)\ge n$. هر ماشین تورینگ چندنواره با زمان $t(n)$ یک ماشین تک‌نواره هم‌ارز با زمان $O(t^2(n))$ دارد.
\end{theorem}
\begin{proofidea}
در فصل ۱۷ دیدیم ماشین تک‌نواره چطور ماشین چندنواره را شبیه‌سازی می‌کند. حالا آن شبیه‌سازی را تحلیل می‌کنیم و نشان می‌دهیم شبیه‌سازی هر گام ماشین چندنواره حداکثر $O(t(n))$ گام طول می‌کشد.
\end{proofidea}
\begin{proof}
فرض کنید $M$ یک ماشین $k$-نواره با زمان $t(n)$ و $S$ شبیه‌ساز تک‌نواره آن (قضیه \ref{thm:multitape}) باشد. $M$ در $t(n)$ گام روی هر نوار حداکثر $t(n)$ خانه را می‌بیند، پس بخش فعال نوار $S$ (که $k$ نوار را کنار هم دارد) طول $O(t(n))$ دارد. در شبیه‌سازی:
\begin{itemize}
\item آماده کردن نوار در ابتدا: $O(n)$ گام.
\item برای هر گام $M$: یک بار طی کردن بخش فعال برای خواندن نمادهای زیر هدها ($O(t(n))$)، یک بار طی کردن برای به‌روزرسانی و حرکت دادن هدها ($O(t(n))$) و حداکثر $k$ بار جابه‌جا کردن محتوا به راست ($O(t(n))$ هر بار). پس هر گام $O(t(n))$.
\end{itemize}
کل شبیه‌سازی $t(n)$ گام از $M$ را شامل می‌شود: $O(n)+t(n)\times O(t(n))=O(t^2(n))$، چون $t(n)\ge n$.
\end{proof}

\begin{theorem}
فرض کنید $t(n)\ge n$. هر ماشین تورینگ غیرقطعی تک‌نواره با زمان $t(n)$ یک ماشین قطعی تک‌نواره هم‌ارز با زمان $2^{O(t(n))}$ دارد.
\end{theorem}
\begin{proof}
فرض کنید $N$ یک ماشین غیرقطعی با زمان $t(n)$ و $D$ ماشین قطعی‌ای باشد که درخت محاسبه $N$ را سطح-به-سطح جست‌وجو می‌کند (قضیه \ref{thm:ntm}). اگر $b$ بیشترین تعداد انتخاب در تابع گذار $N$ باشد، هر گره درخت حداکثر $b$ فرزند دارد و عمق درخت حداکثر $t(n)$ است. پس تعداد کل گره‌ها $O(b^{t(n)})$ است. $D$ برای رسیدن به هر گره، از ریشه شروع می‌کند و حداکثر $t(n)$ گام برمی‌دارد. پس زمان $D$ (که سه نوار دارد) برابر است با
\[O\big(t(n)\,b^{t(n)}\big)=2^{O(t(n))}.\]
تبدیل ماشین سه‌نواره $D$ به تک‌نواره، طبق قضیه قبل زمان را حداکثر به توان دو می‌رساند:
\[\big(2^{O(t(n))}\big)^2=2^{O(2t(n))}=2^{O(t(n))}.\qedhere\]
\end{proof}

\section{کلاس P}

دو قضیه بالا یک تفاوت مهم را نشان می‌دهند: تفاوت زمانی بین ماشین تک‌نواره و چندنواره حداکثر \textbf{چندجمله‌ای} (توان دوم) است، اما تفاوت بین ماشین قطعی و غیرقطعی (تا جایی که می‌دانیم) \textbf{نمایی} است.

\begin{itemize}
\item تفاوت چندجمله‌ای کوچک است؛ تفاوت نمایی بزرگ است.
\item مثلاً برای $n=1000$: $n^3$ یک میلیارد است، اما $2^n$ از تعداد اتم‌های جهان بیشتر است.
\item الگوریتم‌های نمایی به ندرت مفیدند. معمولاً از \textbf{جست‌وجوی کامل} (\lr{brute force}) میان همه جواب‌های ممکن به دست می‌آیند.
\end{itemize}

همه مدل‌های محاسباتی قطعی معقول \textbf{به‌طور چندجمله‌ای هم‌ارز}اند: هر کدام دیگری را با افزایش زمان حداکثر چندجمله‌ای شبیه‌سازی می‌کند. پس روی ویژگی‌هایی از پیچیدگی تمرکز می‌کنیم که با تفاوت‌های چندجمله‌ای عوض نمی‌شوند. به این ترتیب نظریه‌ای داریم که به انتخاب یک مدل خاص وابسته نیست.

\begin{definition}[کلاس P]
$\cP$ کلاس همه زبان‌هایی است که یک ماشین تورینگ قطعی تک‌نواره در \textbf{زمان چندجمله‌ای} تصمیمشان می‌گیرد:
\[\cP=\bigcup_k\mathrm{TIME}(n^k).\]
\end{definition}

\begin{itemize}
\item $\cP$ برای همه مدل‌های محاسباتی‌ای که به‌طور چندجمله‌ای با ماشین تورینگ قطعی تک‌نواره هم‌ارزند، یکسان است. پس $\cP$ از نظر ریاضی کلاسی \textbf{مستحکم} است.
\item $\cP$ تقریباً همان کلاس مسئله‌هایی است که در عمل روی یک کامپیوتر قابل حل‌اند.
\end{itemize}

\begin{theorem}
مسئله مسیر در گراف جهت‌دار در $\cP$ است:
\[\prob{PATH}=\{\enc{G,s,t}\st G\ \text{\rl{گرافی جهت‌دار است که مسیری جهت‌دار از $s$ به $t$ دارد}}\}\in\cP.\]
\end{theorem}
\textbf{جست‌وجوی کامل کافی نیست:} بررسی همه مسیرهای ممکن در گرافی با $m$ رأس کارآمد نیست؛ تعداد مسیرها تقریباً $m^m$ است.
\begin{proof}
الگوریتم چندجمله‌ای (جست‌وجوی سطح-به-سطح):
\begin{algo}{$M$}{روی ورودی $\enc{G,s,t}$ که $G$ گراف جهت‌دار با رأس‌های $s$ و $t$ است}
\begin{enumerate}
\item رأس $s$ را علامت بزن.
\item تا وقتی رأس جدیدی علامت نخورد، تکرار کن:
\item \quad همه یال‌های $G$ را بررسی کن. اگر یالی مثل $(a,b)$ از یک رأس علامت‌خورده $a$ به یک رأس علامت‌نخورده $b$ پیدا شد، $b$ را علامت بزن.
\item اگر $t$ علامت خورده، \textbf{بپذیر}؛ وگرنه \textbf{رد کن}.
\end{enumerate}
\end{algo}
مراحل ۱ و ۴ فقط یک بار اجرا می‌شوند. مرحله ۳ حداکثر $m$ بار اجرا می‌شود (هر بار، جز آخرین، دست‌کم یک رأس علامت می‌خورد). هر مرحله را می‌توان در زمان چندجمله‌ای بر حسب اندازه ورودی پیاده‌سازی کرد. پس $M$ الگوریتمی چندجمله‌ای است.
\end{proof}

\begin{summary}
\begin{itemize}
\item زمان اجرا: بیشترین تعداد گام روی ورودی‌های به طول $n$ (بدترین حالت).
\item $f=O(g)$: $f\le cg$ برای $n$های بزرگ. $f=o(g)$: $f/g\to0$.
\item $\mathrm{TIME}(t(n))$: زبان‌های قابل تصمیم در زمان $O(t(n))$. $\{0^k1^k\}$ در $O(n^2)$، $O(n\log n)$ و با دو نوار در $O(n)$.
\item تبدیل ماشین چندنواره به تک‌نواره زمان را از $t$ به $O(t^2)$ و تبدیل ماشین غیرقطعی به قطعی آن را به $2^{O(t)}$ می‌رساند.
\item $\cP=\bigcup_k\mathrm{TIME}(n^k)$؛ مستحکم و نماینده مسئله‌های عملاً قابل حل. $\prob{PATH}\in\cP$.
\end{itemize}
\end{summary}
